% IoTBench — ACSAC 2026 submission
% Target format: 11 pages double-column IEEEtran conference compsoc
% Double-blind review — NO author info before camera-ready.

\documentclass[conference,compsoc]{IEEEtran}

% ── packages ──────────────────────────────────────────────────────────────
% Using default OT1 encoding — TinyTeX lacks the T1 Courier metric (pcrr8t).
% Once the `psnfss` + `courier` metric files are installed (`tlmgr install courier
% psnfss`), you may re-enable:
%   \usepackage[T1]{fontenc}
\usepackage[utf8]{inputenc}

% TinyTeX also lacks URW Courier OT1 metrics; fall back to Computer Modern
% Typewriter, which is always available.
\renewcommand{\ttdefault}{cmtt}
\usepackage{graphicx}
\usepackage{booktabs}
\usepackage{multirow}
\usepackage{amsmath,amssymb}
\usepackage{xcolor}
\usepackage{xspace}
\usepackage{verbatim}
\usepackage[hidelinks,breaklinks=true]{hyperref}
% Optional polish packages — install with tlmgr once TeX tree is writable:
%   tlmgr install microtype subcaption caption algorithms algorithmicx preprint listings enumitem cleveref

% ── macros ────────────────────────────────────────────────────────────────
\newcommand{\systemname}{IoTBench\xspace}         % benchmark name
\newcommand{\harnessname}{\textsc{IoTPent}\xspace} % harness/agent name (provisional — TBD)
\newcommand{\todo}[1]{\textcolor{red}{\textbf{[TODO: #1]}}}
\newcommand{\note}[1]{\textcolor{blue}{\textit{[note: #1]}}}

% Fallback when cleveref is not installed — plain \ref
\providecommand{\Cref}[1]{Section~\ref{#1}}
\providecommand{\cref}[1]{\S\ref{#1}}

% Remove lstset/lstdefinelanguage: use verbatim for YAML listings instead
\newenvironment{yamlquote}{\par\vspace{.3em}\footnotesize\verbatim}{\endverbatim\normalsize\par\vspace{.3em}}

% ── title & anonymous authors (ACSAC double-blind) ────────────────────────
\title{\systemname: A Network-Scale Benchmark for LLM-Based IoT Penetration Testing}

\author{\IEEEauthorblockN{Anonymous Authors}
\IEEEauthorblockA{Paper~\#XXXX, Anonymized for Double-Blind Review}}

% ──────────────────────────────────────────────────────────────────────────
\begin{document}
\maketitle

\begin{abstract}
% TARGET: 200-250 words
% Structure:
%   1. Context: IoT pentest is slow/manual; LLM agents promise automation.
%   2. Gap: every published LLM-pentest benchmark evaluates isolated hosts.
%   3. Contribution: IoTBench (benchmark) + harness (method) + empirical study.
%   4. Key results with numbers (recall, L2/L3 coverage, cost).
%   5. Artifact: released under OPEN_LICENSE at ANONYMIZED_URL.
\todo{Abstract — write last. Lead with the network-scale gap, not the harness.}
\end{abstract}

\begin{IEEEkeywords}
IoT security, penetration testing, large language models, autonomous agents,
benchmark, cyber-physical systems
\end{IEEEkeywords}

% ── sections ──────────────────────────────────────────────────────────────
% §1 Introduction — target 1 page (~5-6 paragraphs)
% Purpose: plant the network-scale gap and the 4 contributions.
% Do: lead with a concrete IoT-at-scale incident; pose the problem; state the gap; announce contributions.
% Don't: hardware tools, "autonomous exploitation", "graph-guided".

\section{Introduction}
\label{sec:intro}

% ── ¶1 hook: real-world IoT pentest motivation ────────────────────────────
% Concrete incident referencing mass IoT compromise (Mirai variants, smart-city
% camera takeovers, ICS reconnaissance campaigns). Frame as: attacker does not
% compromise *a* device — they compromise *a network* of devices.
\todo{Write ¶1: IoT attacks are network-scale; defenders pentest with single-host tooling.}

% ── ¶2 the promise, the cost ──────────────────────────────────────────────
% Manual IoT pentest is expensive and slow; LLM agents (PentestGPT~\cite{pentestgpt24},
% VulnBot~\cite{vulnbot25}, AutoPentester~\cite{autopentester25}, CAI~\cite{cai25})
% promise automation and have shown strong results on CTF and HackTheBox targets.
\todo{Write ¶2: LLM agents are progressing fast, cite 4-5 recent peers.}

% ── ¶3 the gap ────────────────────────────────────────────────────────────
% Every published LLM-pentest benchmark — AUTOPENBENCH~\cite{autopenbench24},
% HackTheBox, VulnHub, CAIBench/RCTF2~\cite{caibench25} — uses an *isolated host*
% as the unit of evaluation. Real IoT pentest is a network task: heterogeneous
% devices, multiple protocols (MQTT, Modbus, LoRaWAN, HTTP, SSH), lateral reachability
% constrained by architecture. There is no benchmark that evaluates LLM agents at
% this unit.
\todo{Write ¶3: formalize the structural gap. Cite 5-6 benchmarks. Keep it sharp.}

% ── ¶4 contribution statement ─────────────────────────────────────────────
\noindent\textbf{Contributions.} We bridge this gap with:
\begin{itemize}
  \item \textbf{\systemname} — the first reproducible benchmark for LLM pentest
  agents at \emph{network scale}: 5 architectural patterns, 7~scenarios, 8~protocols,
  87~vulnerabilities deterministically injected via Ansible, with ground truth mapped
  to OWASP~IoT~Top~10 and MITRE ATT\&CK for ICS.
  \item \textbf{\harnessname} — a network-native harness with \emph{discovery $\to$
  per-device parallel triage $\to$ per-vulnerability verification $\to$ network-level
  report}, designed to keep LLM context tractable at $/24$ scale.
  \item \textbf{Evidence-level scoring (L1/L2/L3)} — a severity-weighted evaluator
  that distinguishes detection, exploitation, and data exfiltration, beyond the
  binary flag-capture of single-host benchmarks.
  \item \textbf{Empirical study} — using a single shared LLM (\texttt{MiniMax-M2.7})
  to isolate architectural contributions, we compare \harnessname\ against
  three categories of baselines: a non-LLM scanner (Nmap NSE), single-agent
  LLM (PentestGPT), and multi-agent LLM systems with and without graph guidance
  (CAI in two scope variants, VulnBot with task-graph). We additionally report
  two ablations of our pipeline that isolate the individual contributions of
  the infrastructure-graph phase and the multi-hop intrusion phase.
\end{itemize}

% ── ¶5 findings teaser (with numbers) ─────────────────────────────────────
\todo{Write ¶5: one teaser sentence per RQ. e.g. "Best LLM reaches XX\% recall at
/24 scale; scanners reach YY\%; but only ZZ\% of LLM findings achieve L2+."}

% ── ¶6 paper organization ─────────────────────────────────────────────────
\noindent\textbf{Paper organization.}
\Cref{sec:related} surveys LLM pentest agents and benchmarks.
\Cref{sec:threat} formalizes the threat model.
\Cref{sec:iotbench} presents \systemname.
\Cref{sec:harness} describes the harness architecture.
\Cref{sec:setup} and \Cref{sec:eval} present experiments and results.
\Cref{sec:discussion} and \Cref{sec:ethical} discuss limitations and ethics.
Artifacts are released at \texttt{ANONYMIZED\_URL}.

% §2 Background & Related Work — target 1.5 page
% Written 2026-04-23. Positions against 3 clusters.
% Key deliverable: Table 1 comparison grid.

\section{Background and Related Work}
\label{sec:related}

\subsection{LLM-based penetration testing agents}
\label{sec:related:llm}

\paragraph{Single-agent origins}
Deng~\emph{et al.}~proposed \textsc{PentestGPT}~\cite{pentestgpt24}, the first
published LLM-driven penetration testing agent. \textsc{PentestGPT} couples a
single ChatGPT-class model with a human-in-the-loop interface that translates
natural-language guidance into shell commands, and was evaluated on $182$
hand-selected sub-tasks drawn from HackTheBox and VulnHub targets. The work
established three findings that recur throughout the subsequent literature:
LLMs produce plausible reconnaissance plans, they excel on text-rich web
vulnerabilities (SQLi, XSS, IDOR), and they fail on long-context tasks where
state must be maintained across many command invocations.

\paragraph{Multi-agent decomposition}
Later systems address the context limitation through role specialization.
\textsc{VulnBot}~\cite{vulnbot25} decomposes the task into three
phase-specific agents (\emph{reconnaissance}, \emph{scanning}, \emph{exploitation})
coordinated by a Penetration Task Graph (PTG); on the
AUTOPENBENCH corpus~\cite{autopenbench24} its \textsc{Llama3.1-405B} variant
reaches a completion rate of $30.3\%$ versus $9.1\%$ for the same model used
single-agent and $21.2\%$ for GPT-4o.
\textsc{AutoPentester}~\cite{autopentester25} uses five cooperating agents
(summarizer, strategy analyzer, command generator, results verifier, reporter)
with Retrieval-Augmented Generation and a Pentest Tree data structure,
reporting a $27\%$ gain in sub-task completion and $39.5\%$ more vulnerability
coverage over \textsc{PentestGPT} on HackTheBox plus a custom set of OWASP-Top-10
machines.
\textsc{PentestAgent}~\cite{pentestagent25} likewise builds on a multi-stage
pipeline with RAG, published at ACM ASIA~CCS'25.

\paragraph{Learning-based and long-running variants}
\textsc{Pentest-R1}~\cite{pentestr125} applies two-stage reinforcement
learning to the reasoning traces of a mid-scale model, and
\textsc{xOffense}~\cite{xoffense25} fine-tunes a Qwen3-32B with offensive
knowledge corpora, reporting $79\%$ sub-task completion on CTF-style targets.
A complementary strand addresses session duration:
\textsc{CHAP}~\cite{chap26} introduces a context-relay mechanism that
transfers distilled protocols between successive agent instances during
long-running engagements, and \textsc{AWE}~\cite{awe26} adds persistent memory
and browser-backed verification to web-application pentest.
\textsc{CheckMate}~\cite{checkmate25} couples LLM agents with classical
planning to obtain deterministic action selection, improving benchmark
success by over $20$ percentage points.

\paragraph{Production-grade frameworks}
\textsc{CAI}~\cite{cai25}, developed by the \textsc{PentestGPT} authors,
generalizes the single-agent prototype into a bug-bounty-oriented framework
with modular tool support and a public leaderboard; on CAIBench's CTF subset
it consistently outperforms prior open-source systems.

\paragraph{Gap}
Despite this trajectory, every system in this cluster is evaluated on a
\emph{single-host} target unit: HackTheBox machines, VulnHub VMs,
AUTOPENBENCH challenges, or CTF containers. Success is reported either as
flag capture or as sub-task completion within one host.
None evaluates at the scale of a deployed IP network containing heterogeneous
devices and IoT-specific application protocols.

\subsection{Penetration testing benchmarks}
\label{sec:related:bench}

\paragraph{Academic corpora}
AUTOPENBENCH~\cite{autopenbench24} is the current academic reference,
consisting of VulnHub-derived single-VM challenges with scripted flag
checkers; it is used as the canonical evaluation by
\textsc{VulnBot} and \textsc{Pentest-R1}. \textsc{PentestGPT} reported on a
hand-curated set of $182$ sub-tasks drawn from HackTheBox and VulnHub. More
recently, \textsc{CAIBench}~\cite{caibench25} consolidates the field into a
meta-benchmark spanning Jeopardy-style CTFs, attack-and-defense exercises,
cyber-range scenarios, and knowledge/privacy assessments, reporting that
state-of-the-art models saturate around $70\%$ on static knowledge tests but
collapse to $20$--$40\%$ on multi-step adversarial tasks. Within \textsc{CAIBench},
the RCTF2 subset contributes $27$~challenges targeting robotics and
cyber-physical systems --- the first inclusion of CPS in an LLM-pentest
benchmark --- yet each challenge remains an isolated CTF-style target rather
than a deployed network.

\paragraph{Training-oriented lab infrastructures}
\textsc{Ludus}~\cite{ludus} and the \textsc{GOAD}~\cite{goad} Active-Directory
scenario deliver reproducible Proxmox-based labs with Ansible-driven
deployment and idempotent configuration. Methodologically these projects are
the closest precedent to our infrastructure, but they are training tools for
human operators: they do not ship a ground-truth schema, do not define a
scoring procedure, and have not been used to evaluate LLM agents in the
published literature.

\paragraph{Gap}
No benchmark in the published literature provides all of
(i)~network-scale evaluation on a deployed multi-device IP range,
(ii)~IoT-specific application protocols (MQTT, LoRaWAN, Modbus, CoAP),
(iii)~per-vulnerability ground truth usable for automated scoring, and
(iv)~evidence-level granularity beyond binary flag capture. \systemname\ is
designed to fill exactly this intersection.

\subsection{IoT security methodology}
\label{sec:related:iot}

We ground scenario design and ground-truth taxonomy in three widely adopted
frameworks. The OWASP~IoT~Top~10~\cite{owasp-iot} structures our vulnerability
categories and ensures coverage of the attack surfaces most commonly observed
in IoT deployments (hardcoded credentials, insecure network services,
insecure ecosystem interfaces, insufficient update mechanisms). The
MITRE~ATT\&CK~for~ICS~\cite{mitre-ics} matrix provides a tactics vocabulary
for OT-inclusive scenarios; each injected vulnerability in \systemname\ carries
a MITRE-ICS tactic annotation where applicable. Finally,
NIST~SP~800-82~Rev.~3~\cite{nist80082} informs our segmentation and
safety conventions, in particular the IT/OT separation underpinning
architectural pattern~P3 (\cref{sec:iotbench:patterns}).

% ── comparison table (Table 1) ────────────────────────────────────────────
\begin{table*}[t]
\centering
\caption{Comparison with prior LLM penetration testing systems and benchmarks.
``Scope'' is the unit of evaluation; ``Network-scale'' means evaluation
occurs on a multi-device IP network rather than isolated targets.}
\label{tab:related}
\setlength{\tabcolsep}{3pt}
\scriptsize
\begin{tabular}{@{}lccccccc@{}}
\toprule
System & Scope & Protocols & Architectures & IoT-specific & Scoring & Ground truth & Reproducible \\
\midrule
PentestGPT~\cite{pentestgpt24}          & single-host & Web/OS    & n/a       & no      & sub-task      & hand-curated  & partial \\
VulnBot~\cite{vulnbot25}                & single-host & Web/OS    & n/a       & no      & task-graph    & AUTOPENBENCH  & yes     \\
AutoPentester~\cite{autopentester25}    & single-host & Web/OS    & n/a       & no      & sub-task      & HTB + custom  & yes     \\
PentestAgent~\cite{pentestagent25}      & single-host & Web/OS    & n/a       & no      & composite     & custom        & partial \\
Pentest-R1~\cite{pentestr125}           & single-host & Web/OS    & n/a       & no      & sub-task      & AUTOPENBENCH  & partial \\
xOffense~\cite{xoffense25}              & single-host & Web/CTF   & n/a       & no      & sub-task      & CTF           & partial \\
CAI~\cite{cai25}                        & single-host & Web/CTF   & n/a       & no      & CTF           & CAIBench      & yes     \\
CAIBench / RCTF2~\cite{caibench25}      & single-host & mixed     & n/a       & 27 CPS  & CTF           & yes           & yes     \\
Ludus + GOAD~\cite{ludus,goad}          & multi-host  & AD / Win  & AD only   & no      & (training)    & no            & yes     \\
\midrule
\textbf{\systemname\ (ours)}            & \textbf{network} & \textbf{MQTT, Modbus, LoRa, \ldots} & \textbf{5 patterns} & \textbf{yes} & \textbf{L1/L2/L3 + wt.} & \textbf{per-vuln YAML} & \textbf{yes} \\
\bottomrule
\end{tabular}
\end{table*}

% §3 Threat Model & Problem Definition — target 0.5 page
% Purpose: formalize. ACSAC reviewers expect this explicitly.

\section{Threat Model and Problem Definition}
\label{sec:threat}

\paragraph{Network-scale pentest}
We define \emph{network-scale penetration testing} as the task of: (i)~discovering
all live hosts in a target IP range, (ii)~enumerating services and identifying
vulnerabilities on each host, (iii)~producing evidence of exploitability where
feasible, and (iv)~aggregating findings at the network level. The \emph{unit of
evaluation} is the network, not an individual host.

\paragraph{Attacker model}
The attacker holds network-layer access to the target subnet (post-ingress,
pre-compromise). They hold no prior credentials, no privileged vantage, and no
knowledge of deployed software beyond what is observable through scans. The
attacker has access to standard offensive tooling (Nmap, MQTT clients, SSH, \ldots)
and an LLM with tool-calling. Time and budget are bounded.

\paragraph{Goals}
\begin{enumerate}
  \item Enumerate all devices reachable from the entry point.
  \item For each device, identify vulnerabilities from misconfiguration, weak
  credentials, information disclosure, insecure protocols, and known CVEs.
  \item For each finding, attempt to produce \emph{evidence} at one of three
  levels: \textsf{L1}~detection (port open, banner leak), \textsf{L2}~exploitation
  (authenticated session, protocol handshake), \textsf{L3}~data exfiltration
  (configuration files, credentials, sensor data).
  \item Report findings at the network level (severity-ranked, deduplicated).
\end{enumerate}

\paragraph{Scope and non-goals}
We evaluate at the \emph{IP/application layer} on deployed IoT networks. We
do \emph{not} evaluate: (a)~physical/side-channel attacks on hardware,
(b)~RF-layer attacks requiring SDR captures,
(c)~firmware extraction from physical devices,
(d)~post-exploitation persistence or impact,
(e)~insider threats or supply-chain compromise.
Hardware-assisted attacks are left as future work (\cref{sec:conclusion}).

% §4 IoTBench: Benchmark Design — target 2 pages
% Purpose: sell the artifact.
% Key figure: Fig 1 — 5 architectural patterns side by side.
% Key table: Table 2 — scenarios × protocols × vulns × OWASP mapping.

\section{\systemname: A Network-Scale Benchmark}
\label{sec:iotbench}

% ── §4.1 design principles ────────────────────────────────────────────────
\subsection{Design principles}
\label{sec:iotbench:principles}

\systemname{} is built around four principles derived from the gap analysis in
\cref{sec:related}:

\begin{description}
  \item[Reproducibility.] The entire infrastructure deploys from Ansible
  playbooks against a Proxmox cluster; every vulnerability is injected
  idempotently with no manual post-deployment steps.
  \item[Architectural diversity.] Scenarios span five network patterns
  (\cref{fig:patterns}) with distinct attacker reachability graphs enforced by
  OpenWrt firewall rules, not just by service labels.
  \item[Protocol coverage.] Each scenario mixes multiple IoT and OT protocols
  (MQTT, Modbus, LoRaWAN gateways, CoAP, SNMP, plus standard SSH/HTTP/FTP/Telnet)
  to exercise protocol-specific reasoning.
  \item[Ground-truth granularity.] Every injected vulnerability is documented in
  YAML with \texttt{device}, \texttt{IP}, \texttt{role}, \texttt{severity},
  \texttt{category}, \texttt{OWASP IoT} and \texttt{MITRE ICS} mappings,
  optional \texttt{CVE}, \texttt{indicators} (expected detection strings), and
  a \texttt{verification} command.
\end{description}

% ── §4.2 five architectural patterns ──────────────────────────────────────
\subsection{Five architectural patterns}
\label{sec:iotbench:patterns}

\begin{figure*}[t]
  \centering
  % \includegraphics[width=\linewidth]{figures/fig1_patterns.pdf}
  \fbox{\parbox{0.95\linewidth}{\centering\textit{
    Fig.~1 placeholder: five topology diagrams side-by-side.\\
    (Flat)~— (Gateway-centric)~— (Segmented IT/OT)~— (Hub-star)~— (Edge-cloud)
  }}}
  \caption{The five architectural patterns in \systemname. Each pattern is enforced
  by OpenWrt firewall rules that constrain the attacker's reachability from the
  entry point, not by service labels alone.}
  \label{fig:patterns}
\end{figure*}

\todo{Write §4.2 prose — for each of the 5 patterns: (i)~one sentence naming the
pattern, (ii)~the real-world IoT deployment it models, (iii)~the firewall
reachability invariant.}

Concretely:
\begin{enumerate}
  \item \textbf{Flat.} All devices are directly reachable from the entry point;
  no pivot is required.
  \item \textbf{Gateway-centric.} A single IoT gateway mediates all access to
  back-end services; inter-service traffic is blocked without gateway compromise.
  \item \textbf{Segmented IT/OT.} Three zones (admin/IT/OT) with access-control
  lists; the OT zone (Modbus PLCs, HMI) is reachable only via a jump host.
  \item \textbf{Hub-centric star.} A central home-automation hub is the only
  device that can reach MQTT broker, DB, and camera; peripherals cannot talk
  to each other directly.
  \item \textbf{Edge-cloud pivot.} Edge subnet and cloud subnet are separated;
  only the edge gateway bridges them, making cloud-side services unreachable
  without edge compromise.
\end{enumerate}

% ── §4.3 vulnerability injection ──────────────────────────────────────────
\subsection{Deterministic vulnerability injection}
\label{sec:iotbench:inject}

Vulnerabilities are injected by a single Ansible role (\texttt{04\_inject\_vulns.yml})
parameterized by \texttt{scenario\_id}. Each scenario YAML lists \texttt{router\_vulns}
(OpenWrt-level: Telnet, WAN admin, FTP) and per-service \texttt{role} definitions
that trigger role-specific injections: \emph{e.g.}, the \texttt{mqtt\_broker} role
sets \texttt{allow\_anonymous true} and binds on \texttt{0.0.0.0:1883}; the
\texttt{ssh\_server} role installs weak credentials (\texttt{admin:admin}) and
enables \texttt{PasswordAuthentication}; the \texttt{db\_server} role starts
MariaDB with a passwordless \texttt{root}. Category distribution is shown in
\cref{tab:categories}.

\begin{table}[t]
\centering
\caption{Vulnerability categories and severity distribution across \systemname
(87~injected instances).}
\label{tab:categories}
\footnotesize
\begin{tabular}{@{}lrrrr@{}}
\toprule
Category & Count & CRIT & HIGH & MED/LOW \\
\midrule
\texttt{misconfiguration}      & 17 & 3 & 8 & 6 \\
\texttt{data\_exposure}        & 12 & 0 & 7 & 5 \\
\texttt{default\_credentials}  & 11 & 6 & 5 & 0 \\
\texttt{info\_disclosure}      & 10 & 0 & 0 & 10 \\
\texttt{no\_authentication}    &  9 & 2 & 4 & 3 \\
\texttt{insecure\_update}      &  4 & 0 & 0 & 4 \\
\texttt{weak\_crypto}          &  3 & 0 & 0 & 3 \\
\texttt{missing\_header}       &  3 & 0 & 0 & 3 \\
\texttt{known\_cve}            &  3 & 0 & 3 & 0 \\
\texttt{code\_injection}       &  3 & 2 & 1 & 0 \\
\texttt{privilege\_escalation} &  2 & 0 & 2 & 0 \\
\midrule
\textbf{Total}                 &\textbf{87} & 13 & 30 & 44 \\
\bottomrule
\end{tabular}
\end{table}

% ── §4.4 ground truth schema ──────────────────────────────────────────────
\subsection{Ground-truth schema}
\label{sec:iotbench:gt}

Each scenario ships a YAML file describing every injected vulnerability. An
entry example (from \texttt{scenario\_1.yaml}):

\begin{figure}[t]
\footnotesize
\begin{verbatim}
- id: V1
  device: s1-mqtt
  ip: "192.168.100.11"
  role: mqtt_broker
  title: "MQTT without authentication"
  severity: high
  category: misconfiguration
  owasp_iot: "I1 / I9"
  mitre_ics: "Initial Access, Collection"
  cve: null
  indicators:
    - "allow_anonymous true"
    - "Anonymous clients connected"
  verification: "mosquitto_sub -h 192.168.100.11 -t '#' -v"
\end{verbatim}
\caption{Ground-truth entry example.}
\label{lst:gt}
\end{figure}

% ── §4.5 evaluator ────────────────────────────────────────────────────────
\subsection{Evaluator}
\label{sec:iotbench:eval}

The evaluator (\S\,\ref{sec:harness:report} integrates results into Phase~5)
matches each LLM finding against ground truth in four passes, in order:
(1)~exact CVE match; (2)~\texttt{(ip, vuln\_type)}; (3)~\texttt{(ip, category)};
(4)~\texttt{(ip, severity)}. Severity weights are
$w = \{\text{critical}{=}4,\text{high}{=}3,\text{medium}{=}2,\text{low}{=}1\}$;
loose category matches apply $0.5\times$, severity mismatches $0.75\times$.
The composite score is the weighted True-Positive sum normalized by
$\sum_{g \in \text{GT}} w(g)$.

\paragraph{Evidence levels}
Beyond match accuracy, the evaluator tracks, per True Positive, the deepest
\emph{evidence level} achieved by the harness: \textsf{L1}~(detection — banner
or port), \textsf{L2}~(exploitation — authenticated action), \textsf{L3}~(data
exfiltration — sensitive content retrieved). We report
$\text{L\textit{k}-coverage} = |\{t \in \text{TP} : \text{level}(t) \geq k\}| / |\text{TP}|$.
This metric cannot be computed on single-host binary-flag benchmarks and
captures a dimension of pentest quality that existing LLM-pentest evaluations
conflate.

% §5 Harness Architecture — target 2 pages
% Purpose: describe the method. Emphasize: network-native, parallelization, tractable context.
% Key figure: Fig 2 — pipeline 5 phases; Fig 3 — per-device parallelization.

\section{\harnessname: A Network-Native LLM Pentest Harness}
\label{sec:harness}

% ── §5.1 overview ─────────────────────────────────────────────────────────
\subsection{Overview}
\label{sec:harness:overview}

\begin{figure*}[t]
  \centering
  % \includegraphics[width=0.9\linewidth]{figures/fig2_pipeline.pdf}
  \fbox{\parbox{0.9\linewidth}{\centering\textit{
    Fig.~2 placeholder: 5-phase pipeline diagram.\\
    Phase 1 (topology prior) $\to$ Phase 2 (discovery) $\to$ Phase 3 (per-device triage)
    $\to$ Phase 4 (per-vuln verification) $\to$ Phase 5 (network report).
  }}}
  \caption{\harnessname{} pipeline. Solid arrows are deliverables; dashed boxes
  indicate LLM agent calls; Phase~3 and Phase~4 spawn micro-agents in parallel
  (§\ref{sec:harness:phase3}--§\ref{sec:harness:phase4}).}
  \label{fig:pipeline}
\end{figure*}

The harness decomposes network pentest into five phases with typed deliverables.
The key design decision is \emph{parallelization at the device and
vulnerability granularity}: rather than a single monolithic agent that must hold
the full network state, we spawn dedicated micro-agents per target with
constrained tool sets, then aggregate deterministically. This keeps LLM context
bounded at $O(1)$ per device regardless of network size.

% ── §5.2 Phase 1 ──────────────────────────────────────────────────────────
\subsection{Phase~1: Topology prior}
\label{sec:harness:phase1}

Phase~1 loads a graph model of the target network as a directed graph
$G=(V,E)$ where $V$ are devices and $E$ are links labeled by protocol. Two
modes are supported: (i)~\emph{scenario mode} loads the predefined topology
from the benchmark YAML; (ii)~\emph{discovery mode} starts from an empty graph
and instructs Phase~2 to populate $V$ via network scans. The output is a
markdown deliverable enumerating attack surface and preliminary pivot candidates
(betweenness-centrality ranked).

% ── §5.3 Phase 2 ──────────────────────────────────────────────────────────
\subsection{Phase~2: Full-network discovery}
\label{sec:harness:phase2}

Phase~2 executes active reconnaissance across the target \texttt{/24}:
\texttt{nmap\_discovery} (host discovery), \texttt{nmap} (service/version),
\texttt{arp\_scan} (L2 + vendor), and protocol-specific probes
(\texttt{mqtt\_listen} subscribing to \texttt{\#} on discovered brokers,
\texttt{ssh\_audit} against SSH servers, \texttt{curl\_headers} on HTTP).
Unlike single-host harnesses, the agent is not given a target IP — it must
derive one from the CIDR scope.

% ── §5.4 Phase 3 ──────────────────────────────────────────────────────────
\subsection{Phase~3: Per-device parallel triage}
\label{sec:harness:phase3}

\begin{figure}[t]
  \centering
  % \includegraphics[width=\linewidth]{figures/fig3_parallelization.pdf}
  \fbox{\parbox{0.95\linewidth}{\centering\textit{
    Fig.~3 placeholder: Phase 3 parallelization — $n$ devices $\to$ $n$ LLM agents
    in parallel $\to$ deterministic merge.
  }}}
  \caption{Per-device parallelization. Each micro-agent receives only its
  device's scan output, avoiding the context explosion of a monolithic
  network-level agent.}
  \label{fig:parallel}
\end{figure}

Phase~3 runs in three sub-phases:
\begin{enumerate}
  \item \textbf{Deterministic scanner pass.} Ansible-driven subprocess
  wrappers (\texttt{ssh\_audit}, \texttt{curl\_headers}, \texttt{mqtt\_listen},
  \texttt{modbus\_scan}, \ldots) run in parallel per discovered device and
  dump structured scan results to disk.
  \item \textbf{Per-device LLM triage.} For each device, a dedicated LLM agent
  receives only that device's Phase~2+3a outputs and has access to a minimal
  tool set: \texttt{http\_get} (follow-up file retrieval) and
  \texttt{cve\_search} (CPE-to-CVE lookup via NVD). The agent produces a JSON
  vuln list scoped to that device.
  \item \textbf{Deterministic merge.} Per-device outputs are concatenated,
  canonicalized through the taxonomy (\S\ref{sec:harness:taxonomy}), and
  deduplicated via severity-aware rules.
\end{enumerate}

This decomposition is essential for network scale: a monolithic agent on a
15-device scenario would need to fit $\sim$30--80k tokens of scan output into
a single context and reason about cross-device relationships
simultaneously. Per-device micro-agents flatten this to a constant size, at
the cost of losing some cross-device inference capability, which is recovered
in Phase~5.

% ── §5.5 Phase 4 ──────────────────────────────────────────────────────────
\subsection{Phase~4: Per-vulnerability verification with evidence levels}
\label{sec:harness:phase4}

Phase~4 spawns one micro-agent per Phase~3 finding (in parallel, with a
concurrency cap). Each micro-agent receives the finding (IP, type, port,
details) and full operational tool access: \texttt{ssh\_login},
\texttt{mysql\_query}, \texttt{mqtt\_listen}, \texttt{http\_get},
\texttt{ftp\_list}, \texttt{modbus\_scan}, \texttt{telnet\_connect}. The agent
is prompted to produce the deepest evidence level it can (L1--L3,
\S\ref{sec:iotbench:eval}) within a bounded turn budget.
Outputs include: \texttt{status} (\textsc{confirmed}, \textsc{not\_exploitable},
\textsc{error}), \texttt{evidence\_level}, \texttt{command} executed,
\texttt{stdout} excerpts, and \texttt{data\_extracted} (when L3).

Findings in \textsc{confirmed} state from Phase~3 override Phase~4
\textsc{error}/\textsc{failed} statuses (rationale:
Phase~3 is hypothesis generation with directly-observed indicators; Phase~4
crashes should not erase directly-observed findings).

% ── §5.6 Phase 5 ──────────────────────────────────────────────────────────
\subsection{Phase~5: Network-level report}
\label{sec:harness:report}

Phase~5 consumes all prior deliverables and produces a network-scoped report
with: device inventory, per-device vulnerabilities with severity, attack
surface summary, and actionable recommendations ordered by severity $\times$
exposure. This is the only phase with cross-device reasoning, deliberately
placed at the end to avoid bloating earlier phases.

% ── §5.7 taxonomy & tool abstraction ──────────────────────────────────────
\subsection{Taxonomy and tool abstraction}
\label{sec:harness:taxonomy}

A single module resolves LLM-produced vulnerability type strings into a
canonical taxonomy of 11~categories (\cref{tab:categories}), filtering noise
(e.g., meta-observations) and CONFIG-only findings. Tools are defined
declaratively as YAML (schema: name, description, JSON-schema parameters,
subprocess command template), so adding a new protocol probe requires no
Python code.

% ── §5.8 complexity ───────────────────────────────────────────────────────
\subsection{Complexity and cost envelope}
\label{sec:harness:cost}

For a network with $n$ devices and an average of $v$ findings per device, the
harness makes $O(1+n+nv)$ LLM calls (Phases 1, 3b, 4 respectively), bounded
per call by fixed turn budgets ($50/50/10/10/20$ for phases 1--5). Total cost
scales linearly with network size, not combinatorially.

% §6 Experimental Setup — target 0.75 page

\section{Experimental Setup}
\label{sec:setup}

\paragraph{Infrastructure}
All scenarios run on a single Proxmox VE host. Networks are Linux bridges
(\texttt{vmbr1}, \texttt{192.168.100.0/24}) with an OpenWrt router as gateway.
Per-scenario deployment and teardown are automated via Ansible and complete
in under 10~minutes per scenario.

\paragraph{LLM model — fairness lock}
To isolate \emph{architectural} contributions from model-specific advantages, all
LLM-driven systems compared in this work --- our pipeline (and its ablations) and
all LLM baselines --- run on the \emph{same} general-purpose model:
\textbf{MiniMax-M2.7} \cite{minimax27}. We deliberately avoid a multi-model
matrix: a fair comparison of architectures requires the model variable to be
held constant. CAI's specialised \texttt{alias1} model (security-fine-tuned,
paywalled) is intentionally not used --- this study addresses the architecture
axis, not the model fine-tuning axis. The choice of MiniMax-M2.7 also keeps the
combined API cost at zero under our research subscription, enabling the full
$k{=}3$-run variance estimation across all systems and scenarios without budget
trade-offs.

\paragraph{Baselines}
We compare our pipeline against three categories of systems, each occupying a
distinct architectural slot:

\begin{itemize}
  \item \textbf{Single-agent LLM} --- \textsc{PentestGPT}~\cite{pentestgpt24}
  in autonomous mode, the canonical USENIX'24 reference.
  \item \textbf{Multi-agent LLM (no infrastructure graph)} ---
  \textsc{CAI}~\cite{cai25}. We test two variants: \emph{A1} per-IP
  (one independent CAI session per ground-truth IP, no shared state) and
  \emph{B} CIDR-scope (a single CAI session given the full \texttt{/24}, allowed
  to discover hosts and pivot via \texttt{LinuxCmd} primitives). Both variants
  receive the same total turn budget as our pipeline.
  \item \textbf{Multi-agent LLM with task-graph} ---
  \textsc{VulnBot}~\cite{vulnbot25}, the most structurally similar competitor.
  VulnBot's Penetration Task Graph (PTG) is a directed acyclic graph of
  \emph{tasks} (action dependencies), to be contrasted with our infrastructure
  graph of \emph{devices} (firewall-enforced reachability). Both are
  ``graph-based'', but encode orthogonal abstractions.
  \item \textbf{Non-LLM scanner} --- Nmap with all NSE scripts
  (\texttt{--script=default,vuln,auth}) against the same target CIDR. Scanner
  outputs are converted to the evaluator's JSON schema.
\end{itemize}

\paragraph{Internal ablations}
Two ablations isolate the contribution of our two key design choices:
(i)~\harnessname--w/o-Phase~1 disables graph analysis (\texttt{--phases 2 3 4 5
6}), testing the value of infrastructure-graph guidance;
(ii)~\harnessname--w/o-Phase~5 disables the intrusion phase (\texttt{--phases 1
2 3 4 6}), testing the value of structured multi-hop lateral movement.

\paragraph{Variance}
Each (system, scenario) cell is run $k{=}3$ times. We report mean, standard
deviation, and 95\% bootstrap confidence intervals. Pairwise differences
(our pipeline vs each baseline) are tested with Mann--Whitney $U$ on F1 and
weighted Score per scenario. This protocol exceeds the single-run reporting of
PentestGPT~\cite{pentestgpt24}, AutoPentester~\cite{autopentester25},
AUTOPENBENCH~\cite{autopenbench24}, and CAI~\cite{cai25}, and matches the
multi-run averaging of VulnBot~\cite{vulnbot25} (5-experiment) and
Pentest-R1~\cite{pentestr125} (pass@3).

\paragraph{Metrics}
We report five families of metrics, chosen to (i)~align with prior work's
expected vocabulary and (ii)~surface architectural differences invisible to
host-level benchmarks:
\begin{itemize}
  \item \emph{Task completion} --- Recall against the ground-truth vulnerability
  set, the network-scale analogue of the sub-task / task completion rate
  reported by PentestGPT, VulnBot, AutoPentester, Pentest-R1, and xOffense.
  \item \emph{Quality} --- Precision, F1, and weighted Score
  (CVSS-severity-weighted Recall, \S\ref{sec:iotbench:eval}). Score plays the
  role of AutoPentester's ``vulnerability coverage'' but with severity weights.
  \item \emph{Multi-Hop Reach} --- MHR$_k$ for $k\in\{1,2,3\}$, the fraction of
  ground-truth vulnerabilities at \texttt{hop\_depth}~$\geq k$ that are reached.
  No prior LLM-pentest benchmark measures network depth explicitly.
  \item \emph{Evidence depth} --- $L_1$ (detection), $L_2$ (exploitation), $L_3$
  (exfiltration) coverage; absent from binary flag-capture benchmarks.
  \item \emph{Cost envelope} --- prompt+completion tokens and monetary cost
  (USD) from provider billing APIs, plus wall-clock time per scenario,
  matching CAI~\cite{cai25}, AutoPentester, and Pentest-R1.
\end{itemize}

\paragraph{Compute and budget}
All runs consume a combined API budget under \$NNN, bounded by per-phase turn
caps and the concurrency limit in Phase~4.
\todo{Insert exact budget after all runs complete (freeze 2 May).}

% §7 Evaluation — target 2.25 pages
% Structure mirrors RQ1-RQ4.

\section{Evaluation}
\label{sec:eval}

% ── §7.1 overall ──────────────────────────────────────────────────────────
\subsection{Overall performance (RQ1)}
\label{sec:eval:overall}

\begin{table*}[t]
\centering
\caption{Main results — Recall / Precision / F1 / weighted Score, averaged over
7 scenarios $\times$ 3 runs per model. Weighted Score normalizes by
$\sum w(g)$ for $g \in \text{GT}$.}
\label{tab:main}
\footnotesize
\begin{tabular}{@{}lrrrrr@{}}
\toprule
Model & Recall & Precision & F1 & Score & Cost (\$) \\
\midrule
Claude Sonnet 4     & \todo{} & \todo{} & \todo{} & \todo{} & \todo{} \\
Gemini 2.5          & \todo{} & \todo{} & \todo{} & \todo{} & \todo{} \\
GPT-4.1             & \todo{} & \todo{} & \todo{} & \todo{} & \todo{} \\
DeepSeek-V3         & \todo{} & \todo{} & \todo{} & \todo{} & \todo{} \\
MiniMax/Qwen3       & \todo{} & \todo{} & \todo{} & \todo{} & \todo{} \\
\bottomrule
\end{tabular}
\end{table*}

\todo{Write §7.1 prose — headline finding: best LLM reaches XX\% recall at /24
scale, variance across runs stays under YY\%.}

% ── §7.2 baselines ────────────────────────────────────────────────────────
\subsection{Network-scale vs baselines (RQ4)}
\label{sec:eval:baselines}

\begin{figure}[t]
  \centering
  % \includegraphics[width=\linewidth]{figures/fig4_baselines.pdf}
  \fbox{\parbox{0.95\linewidth}{\centering\textit{Fig.~4 placeholder: bar chart,
  F1 and Score for \harnessname vs OpenVAS, Nmap NSE, CAI-adapted,
  VulnBot-adapted.}}}
  \caption{Baseline comparison. Scanners capture known CVEs but miss
  configuration vulnerabilities. Single-host LLM agents lose cross-device
  signal when adapted to a network scope.}
  \label{fig:baselines}
\end{figure}

\todo{Write §7.2 — explicit narrative: (i) scanners win on CVE category,
(ii) LLMs win on misconfig/default-creds, (iii) single-host LLMs adapted to
network fall behind native harness on F1 by ZZ\%.}

% ── §7.3 per-architecture ─────────────────────────────────────────────────
\subsection{Per-architecture analysis (RQ3)}
\label{sec:eval:arch}

\begin{figure}[t]
  \centering
  % \includegraphics[width=\linewidth]{figures/fig5_heatmap.pdf}
  \fbox{\parbox{0.95\linewidth}{\centering\textit{Fig.~5 placeholder:
  heatmap scenario $\times$ model, cell = weighted Score.}}}
  \caption{Per-scenario Score heatmap. \todo{Note which pattern consistently
  breaks LLMs.}}
  \label{fig:heatmap}
\end{figure}

\todo{Write §7.3 — pattern-by-pattern commentary. Expected finding:
gateway-centric and hub-star degrade most, because they require pivot
reasoning. Edge-cloud may also degrade on cloud-side service discovery.}

% ── §7.4 evidence level depth ─────────────────────────────────────────────
\subsection{Evidence-level depth (RQ2)}
\label{sec:eval:evidence}

\begin{table}[t]
\centering
\caption{Evidence-level coverage. $L_k$ is the fraction of true positives
reaching level~$k$ or deeper. Absent from binary-flag benchmarks.}
\label{tab:evidence}
\footnotesize
\begin{tabular}{@{}lrrr@{}}
\toprule
Model & $L_1$ (detect) & $L_2$ (exploit) & $L_3$ (exfiltrate) \\
\midrule
Claude Sonnet 4  & \todo{} & \todo{} & \todo{} \\
Gemini 2.5       & \todo{} & \todo{} & \todo{} \\
GPT-4.1          & \todo{} & \todo{} & \todo{} \\
DeepSeek-V3      & \todo{} & \todo{} & \todo{} \\
MiniMax/Qwen3    & \todo{} & \todo{} & \todo{} \\
\midrule
OpenVAS          & --      & --      & --      \\
Nmap NSE         & --      & --      & --      \\
\bottomrule
\end{tabular}
\end{table}

\todo{Write §7.4 — expected: detection (L1) near-parity with scanners, but
L2/L3 is LLM-only. This is the strongest argument for LLM agents in this
context.}

% ── §7.5 per-protocol ─────────────────────────────────────────────────────
\subsection{Per-protocol analysis}
\label{sec:eval:protocol}

\todo{Write §7.5 — 2 paragraphs. Expected: LLMs excel on MQTT/SSH/HTTP
(text-rich, well-represented in training), fall short on Modbus/LoRaWAN
(binary, sparse training data). Report numbers per protocol.}

% ── §7.6 failure modes, cost, time ────────────────────────────────────────
\subsection{Failure modes and cost/time envelope}
\label{sec:eval:failures}

\begin{figure}[t]
  \centering
  % \includegraphics[width=\linewidth]{figures/fig6_pareto.pdf}
  \fbox{\parbox{0.95\linewidth}{\centering\textit{Fig.~6 placeholder:
  Pareto plot, USD cost vs weighted Score, one point per model.}}}
  \caption{Cost--performance Pareto. \todo{Annotate which model dominates
  on each axis and the efficient frontier.}}
  \label{fig:pareto}
\end{figure}

\todo{Write §7.6 — categorize failure modes: (a) hallucinated findings on
devices with no services, (b) missed vulnerabilities on binary protocols,
(c) over-reporting info-disclosure without severity calibration. Brief
quantitative table.}

% §8 Discussion & Limitations — target 0.5 page

\section{Discussion and Limitations}
\label{sec:discussion}

\paragraph{Where LLMs help, where they don't}
\systemname{} surfaces an empirical pattern: LLM agents reach high recall on
misconfigurations, default credentials, and information disclosure — findings
that are textually expressible and well represented in their training corpora.
They consistently underperform on binary industrial protocols (Modbus, CoAP,
LoRaWAN) and on cross-device inferences that require holding an aggregated
network view. This suggests a hybrid deployment model: LLM agents for
configuration auditing, traditional scanners for CVE coverage, and human
operators for the OT/ICS layer.

\paragraph{Implications for agentic systems}
The harness's per-device and per-vulnerability micro-agent structure is one
answer to context scaling, but it cedes cross-device inference. Designs that
keep a global aggregator agent in the loop without paying the full
network-context cost are an open research direction.

\paragraph{Limitations}
(i)~\emph{Simulated infrastructure.} \systemname{} scenarios emulate IoT
devices on Linux VMs; real constrained hardware introduces embedded-specific
attack surfaces (bootloaders, RTOS, firmware) we do not cover.
(ii)~\emph{Evaluator leniency.} Bonus categories are accepted as non-FP even
without ground-truth matches, which inflates precision.
(iii)~\emph{Model variance.} With $k{=}3$ runs per configuration, confidence
intervals on F1 remain wide for some scenarios.
(iv)~\emph{Prompt sensitivity.} We fix a single prompt family across models;
per-model prompt tuning would likely raise scores.
(v)~\emph{Static architectures.} The five patterns are fixed; real deployments
evolve over time and include devices not represented here (RFID badge readers,
industrial PLCs, medical devices).

% §9 Ethical Considerations — target 0.25 page (mandatory per ACSAC CFP)

\section{Ethical Considerations}
\label{sec:ethical}

\paragraph{Isolation}
All experiments run on an isolated Proxmox cluster with a dedicated bridge
(\texttt{vmbr1}) and no route to external networks; the OpenWrt gateway's
\texttt{WAN} interface is firewall-isolated. No live hosts, production
services, or third-party infrastructure are probed.

\paragraph{Responsible disclosure}
\systemname{} vulnerabilities are deterministically injected into our own
lab VMs via Ansible. No 0-days are introduced; all categories (misconfig,
weak credentials, Terrapin CVE-2023-48795, outdated Dropbear) correspond to
patterns already publicly documented. No responsible-disclosure process is
triggered.

\paragraph{Misuse risk of released artifact}
The benchmark artifact (Ansible playbooks, ground truth, evaluator) does not
embed exploitation payloads for unknown vulnerabilities, and the injected
weaknesses are common misconfigurations documented in OWASP~IoT~Top~10.
The harness uses only standard pentest tools already in widespread use.
We judge the net societal benefit — reproducible evaluation of LLM pentest
agents, and a concrete evaluation surface for future defenders — to exceed
the incremental misuse risk.

\paragraph{IRB / human subjects}
No human subjects are involved. No personal data is collected or processed.

\paragraph{LLM costs and environmental footprint}
All reported experiments together consume an estimated \todo{X} kWh and
cost under \todo{USD Y} in API billing.

% §10 Conclusion — target 0.25 page

\section{Conclusion}
\label{sec:conclusion}

We presented \systemname, the first reproducible benchmark for evaluating LLM
penetration testing agents at \emph{network scale} on IoT infrastructure, and
\harnessname, a network-native harness that keeps LLM context tractable
through per-device and per-vulnerability parallelization. Our empirical study
of five state-of-the-art models across 7~scenarios and three architectural
patterns\todo{(update to five once firewall rules land)} shows that LLM agents
outperform traditional scanners on configuration and credential categories
while underperforming on industrial protocols, and that the
detection-to-exfiltration (L1--L3) evidence spectrum exposes model
differences that binary-flag benchmarks conflate.

\paragraph{Future work}
(i)~\emph{Credential propagation} — closing the Phase~4 $\to$ Phase~2 loop to
enable multi-hop exploitation chains;
(ii)~\emph{Real hardware} — extending to embedded Zigbee, LoRaWAN, and
sub-GHz attack surfaces via SDR tooling;
(iii)~\emph{Defender's side} — using \systemname{} as a benchmark for
LLM-assisted IoT defense (detection, patching recommendation).

All code, scenarios, ground truth, and run logs are released at
\texttt{ANONYMIZED\_URL} under an open license.


% ── bibliography ──────────────────────────────────────────────────────────
\bibliographystyle{IEEEtran}
\bibliography{references}

% ── LLM Usage Statement (outside page quota per ACSAC CFP) ────────────────
% §11 LLM Usage Statement — outside page-limit per ACSAC CFP
% The statement must cover BOTH roles: LLM as subject of study AND LLM as writing aid.

\section*{LLM Usage Statement}
\label{sec:llm-usage}

This paper involves Large Language Models in two distinct roles.

\paragraph{LLM as subject of study}
The experimental work evaluates the penetration testing capabilities of LLM
agents using \texttt{MiniMax-M2.7}~\cite{minimax27} as the shared model across
\harnessname\ and all baselines (CAI, VulnBot, PentestGPT) to isolate the
architectural contribution from the model variable (see \S\ref{sec:setup}).
All inference calls were made through MiniMax's official API under standard
terms of service. No model weights were fine-tuned; all systems were used
out-of-the-box with identical tool schemas and prompt templates per system.
Prompt templates, full conversation transcripts, and per-phase token counts
are released as part of the artifact.

\paragraph{LLM as writing aid}
During manuscript preparation, the authors used current frontier LLMs
(Anthropic Claude Opus 4.7, OpenAI GPT-5.5) for: (i)~drafting candidate
section outlines subsequently edited by the authors; (ii)~copy-editing for
grammar, clarity, and concision; (iii)~consistency checks on terminology and
notation. No section was written end-to-end by an LLM. All claims,
experimental results, related-work positioning, and scientific interpretations
are the authors' own. Every LLM-suggested citation was verified by retrieving
the cited work.


\end{document}
